
\begin{center}
{\LARGE\bfseries Arbeitsblatt zum Lernpfad: Excel-Tabellen}\\[1mm]
{\large Klasse 7 -- Schulfest-Kiosk planen und auswerten}
\end{center}

\begin{boxK}
\textbf{Name:} \linie[5.8cm] \hfill \textbf{Datum:} \linie[4.4cm]

\textbf{So nutzt du dieses Arbeitsblatt:} Bearbeite die Etappen im digitalen Lernpfad nacheinander. Trage deine Ergebnisse hier ein. Wenn du bei einer Aufgabe nicht weiterkommst, nutzt du zuerst die \textbf{Excel-Hilfe} auf dieser Seite. Danach prüfst du deine Lösung mit der Kontrollfrage und erst dann mit der Kontrolllösung im Lernpfad.
\end{boxK}

\section*{Kurze Excel-Hilfe}
\begin{boxHinweis}
\textbf{Diese Übersicht ist deine Hilfe für alle Etappen.} Schaue immer wieder hier nach, bevor du rätst oder nur eine Formel abschreibst.
\end{boxHinweis}

\subsection*{1. Tabelle lesen}
\begin{tabularx}{\textwidth}{|L{3.0cm}|Y|L{3.0cm}|}
\hline
\rowcolor{gray!20}\textbf{Begriff} & \textbf{Bedeutung} & \textbf{Beispiel}\\\hline
Arbeitsmappe & die ganze Excel-Datei & Schulfest-Kiosk.xlsx\\\hline
Tabellenblatt & eine einzelne Seite in der Datei & Messreihe, Schueler\_Vorlage\\\hline
Spalte & senkrechter Bereich & A, B, C\\\hline
Zeile & waagerechter Bereich & 1, 2, 3\\\hline
Zelle & Schnittpunkt aus Spalte und Zeile & C3\\\hline
Zellbezug & Verweis auf eine bestimmte Zelle & C3 bedeutet: Wert aus C3\\\hline
\end{tabularx}

\subsection*{2. Daten richtig eintragen}
\begin{boxK}
\textbf{Regel:} Lies zuerst die \textbf{Zeile}, danach die \textbf{Spalte}. Trage den Wert dort ein, wo sich Zeile und Spalte treffen.

\textbf{Beispiel:} Der Messwert für \textbf{Zeile 1, Spalte 3} lautet \textbf{12}. Dann gehört die Zahl 12 in das Feld bei Zeile 1 und Spalte 3.
\end{boxK}

\subsection*{3. Formeln und Befehle verstehen}
\begin{tabularx}{\textwidth}{|L{3.1cm}|L{4.1cm}|Y|}
\hline
\rowcolor{gray!20}\textbf{Zeichen/Befehl} & \textbf{Beispiel} & \textbf{Bedeutung}\\\hline
\code{=} & \code{=C3*D3} & Excel soll rechnen. Ohne \code{=} bleibt es Text.\\\hline
\code{*} & \code{=C3*D3} & Wert aus C3 mal Wert aus D3.\\\hline
\code{-} & \code{=E3-F3} & Wert aus F3 wird von E3 abgezogen.\\\hline
\code{:} & \code{G3:G8} & alle Zellen von G3 bis G8.\\\hline
\code{SUMME} & \code{=SUMME(G3:G8)} & alle Werte von G3 bis G8 addieren.\\\hline
\code{MITTELWERT} & \code{=MITTELWERT(G3:G8)} & Durchschnitt der Werte berechnen.\\\hline
\texttt{\$} & \texttt{\$B\$13} & Diese Zelle bleibt beim Herunterziehen fest.\\\hline
\end{tabularx}

\begin{boxTipp}
\textbf{Drei Fragen vor jeder Formel:}
\begin{enumerate}[label=\arabic*)]
\item Welche Werte brauche ich?
\item In welchen Zellen stehen diese Werte?
\item Welche Rechenoperation passt dazu?
\end{enumerate}
\end{boxTipp}

\medskip
\section*{Etappe 0: Orientierung}
\begin{boxHinweis}
\textbf{Aufgabe.} Öffne die Excel-Datei. Lies die Überschriften. Notiere, was mit der Tabelle herausgefunden werden soll.
\end{boxHinweis}

\textbf{Mein Ziel der Tabelle:}\\[0.8cm]
\linie\\[0.9cm]
\linie

\begin{boxTipp}
\textbf{Hilfe-Verweis:} Wenn du unsicher bist, schaue in der Excel-Hilfe bei \textbf{Tabelle lesen} nach.
\end{boxTipp}

\textbf{Kontrollfrage:} Was ist in dieser Stunde wichtiger: der höchste Verkaufspreis oder der größte Gewinn? Begründe kurz.\\[1cm]
\linie\\[0.8cm]
\linie

\section*{Etappe 1: Excel-Grundlagen}
\begin{boxHinweis}
\textbf{Aufgabe.} Ergänze die Beispiele und erkläre die Befehle in eigenen Worten.
\end{boxHinweis}

\begin{tabularx}{\textwidth}{|L{3.1cm}|Y|L{3.5cm}|}
\hline
\rowcolor{gray!20}\textbf{Begriff} & \textbf{Bedeutung} & \textbf{Mein Beispiel}\\\hline
Spalte & senkrechter Bereich & \antwortfeld[2.8cm]\\\hline
Zeile & waagerechter Bereich & \antwortfeld[2.8cm]\\\hline
Zelle & Schnittpunkt aus Spalte und Zeile & \antwortfeld[2.8cm]\\\hline
Zellbezug & Verweis auf eine Zelle & \antwortfeld[2.8cm]\\\hline
Formel & Rechenbefehl für Excel & \antwortfeld[2.8cm]\\\hline
\end{tabularx}

\vspace{0.2cm}
\textbf{Erkläre kurz:} Warum beginnt eine Formel mit \code{=}?\\[0.8cm]
\linie\\[0.8cm]
\linie

\textbf{Kontrollfragen:} 
\begin{enumerate}[label=\alph*)]
\item Die Zelle in Spalte C und Zeile 3 heißt: \antwortfeld[2cm]
\item Excel nutzt für ``mal'' das Zeichen: \antwortfeld[2cm]
\item Der Bereich \code{G3:G8} bedeutet: \linie[7cm]
\end{enumerate}

\medskip
\section*{Etappe 2: Messwerte in eine Tabelle eintragen}
\begin{boxHinweis}
\textbf{Aufgabe.} Trage die Messwerte aus der Liste in die leere Tabelle ein. Es wird noch nichts berechnet.
\end{boxHinweis}

\begin{minipage}[t]{0.45\textwidth}
\textbf{Messwertliste}
\begin{tabularx}{\textwidth}{|Y|L{1.8cm}|}
\hline
\rowcolor{gray!20}\textbf{Angabe} & \textbf{Wert}\\\hline
Zeile 1, Spalte 1 & 8\\\hline
Zeile 1, Spalte 2 & 10\\\hline
Zeile 1, Spalte 3 & 12\\\hline
Zeile 2, Spalte 1 & 7\\\hline
Zeile 2, Spalte 2 & 11\\\hline
Zeile 2, Spalte 3 & 15\\\hline
Zeile 3, Spalte 1 & 9\\\hline
Zeile 3, Spalte 2 & 13\\\hline
Zeile 3, Spalte 3 & 14\\\hline
\end{tabularx}
\end{minipage}\hfill
\begin{minipage}[t]{0.50\textwidth}
\textbf{Meine Tabelle}
\begin{tabular}{|c|c|c|c|}
\hline
 & \textbf{Spalte 1} & \textbf{Spalte 2} & \textbf{Spalte 3}\\\hline
\textbf{Zeile 1} & \minifeld & \minifeld & \minifeld\\\hline
\textbf{Zeile 2} & \minifeld & \minifeld & \minifeld\\\hline
\textbf{Zeile 3} & \minifeld & \minifeld & \minifeld\\\hline
\end{tabular}
\end{minipage}

\begin{boxTipp}
\textbf{Hilfe-Verweis:} Nutze die Excel-Hilfe bei \textbf{Daten richtig eintragen}: erst Zeile, dann Spalte.
\end{boxTipp}

\textbf{Kontrollfragen:}
\begin{enumerate}[label=\alph*)]
\item Welcher Wert steht bei Zeile 2, Spalte 3? \antwortfeld[2cm]
\item Welcher Wert steht bei Zeile 3, Spalte 2? \antwortfeld[2cm]
\item Wo steht der Wert 12? \linie[6cm]
\end{enumerate}

\section*{Etappe 3: Schulfest-Kiosk-Tabelle lesen}
\begin{boxHinweis}
\textbf{Aufgabe.} Erkläre die Größen in der Kiosk-Tabelle.
\end{boxHinweis}
\begin{tabularx}{\textwidth}{|L{4cm}|Y|}
\hline
\rowcolor{gray!20}\textbf{Größe} & \textbf{Meine Erklärung}\\\hline
Einkaufspreis & \\\hline
Verkaufspreis & \\\hline
Verkaufte Anzahl & \\\hline
Umsatz & \\\hline
Kosten & \\\hline
Gewinn & \\\hline
\end{tabularx}

\begin{boxTipp}
\textbf{Hilfe-Verweis:} Bei Spalten und Zellbezügen hilft dir die Excel-Hilfe unter \textbf{Tabelle lesen}.
\end{boxTipp}

\textbf{Kontrollfragen:}
\begin{enumerate}[label=\alph*)]
\item In welcher Spalte steht der Verkaufspreis? \antwortfeld[2cm]
\item Welche zwei Größen brauchst du für den Umsatz? \linie[7cm]
\end{enumerate}

\medskip
\section*{Etappe 4: Umsatz berechnen}
\begin{boxHinweis}
\textbf{Aufgabe.} Berechne in Excel den Umsatz. Erkläre vorher die Formel.
\end{boxHinweis}

\begin{tabularx}{\textwidth}{|L{5cm}|Y|}
\hline
\rowcolor{gray!20}\textbf{Frage vor der Formel} & \textbf{Meine Antwort}\\\hline
Welche Werte brauche ich? & \\\hline
In welchen Zellen stehen sie beim Wasser? & \\\hline
Welche Rechnung passt? & \\\hline
Welche Formel kommt in E3? & \\\hline
\end{tabularx}

\begin{boxTipp}
\textbf{Hilfe-Verweis:} In der Excel-Hilfe findest du \code{=} und \code{*}. Formel erst eintragen, wenn du C3 und D3 erklären kannst.
\end{boxTipp}

\textbf{Meine Umsatzwerte:}
\begin{tabularx}{\textwidth}{|Y|L{3cm}|Y|L{3cm}|}
\hline
Wasser & \antwortfeld[2cm] & Brezel & \antwortfeld[2cm]\\\hline
Apfelschorle & \antwortfeld[2cm] & Muffin & \antwortfeld[2cm]\\\hline
Eistee & \antwortfeld[2cm] & Obstbecher & \antwortfeld[2cm]\\\hline
\end{tabularx}

\textbf{Kontrollfrage:} Welcher Umsatz entsteht bei Wasser? \antwortfeld[2.5cm]

\section*{Etappe 5: Kosten und Gewinn berechnen}
\begin{boxHinweis}
\textbf{Aufgabe.} Berechne Kosten und Gewinn. Erkläre beide Formeln.
\end{boxHinweis}

\begin{tabularx}{\textwidth}{|L{4.1cm}|Y|}
\hline
\rowcolor{gray!20}\textbf{Formel} & \textbf{Erklärung}\\\hline
Kosten in F3: \code{=B3*D3} & \\\hline
Gewinn in G3: \code{=E3-F3} & \\\hline
\end{tabularx}

\begin{boxTipp}
\textbf{Hilfe-Verweis:} Für Kosten brauchst du \code{*}; für Gewinn brauchst du \code{-}. Beides steht in der Excel-Hilfe.
\end{boxTipp}

\textbf{Meine Ergebnisse:}
\begin{tabularx}{\textwidth}{|Y|L{2.5cm}|L{2.5cm}|}
\hline
\rowcolor{gray!20}\textbf{Artikel} & \textbf{Kosten} & \textbf{Gewinn}\\\hline
Wasser & & \\\hline
Apfelschorle & & \\\hline
Eistee & & \\\hline
Brezel & & \\\hline
Muffin & & \\\hline
Obstbecher & & \\\hline
\end{tabularx}

\textbf{Kontrollfragen:}
\begin{enumerate}[label=\alph*)]
\item Gewinn Wasser: \antwortfeld[2.5cm]
\item Gewinn Brezel: \antwortfeld[2.5cm]
\end{enumerate}

\medskip
\section*{Etappe 6: Summen und Mittelwert}
\begin{boxHinweis}
\textbf{Aufgabe.} Berechne die Summen in Zeile 10 und den durchschnittlichen Gewinn.
\end{boxHinweis}

\begin{tabularx}{\textwidth}{|L{3.5cm}|L{4.2cm}|Y|}
\hline
\rowcolor{gray!20}\textbf{Zelle} & \textbf{Formel} & \textbf{Mein Ergebnis}\\\hline
D10 & \code{=SUMME(D3:D8)} & \\\hline
E10 & \code{=SUMME(E3:E8)} & \\\hline
F10 & \code{=SUMME(F3:F8)} & \\\hline
G10 & \code{=SUMME(G3:G8)} & \\\hline
B12 & \code{=MITTELWERT(G3:G8)} & \\\hline
\end{tabularx}

\begin{boxTipp}
\textbf{Hilfe-Verweis:} Wenn du nicht weißt, was \code{G3:G8}, \code{SUMME} oder \code{MITTELWERT} bedeutet, schaue in die Excel-Hilfe.
\end{boxTipp}

\textbf{Kontrollfragen:}
\begin{enumerate}[label=\alph*)]
\item Gesamtgewinn: \antwortfeld[2.5cm]
\item Durchschnittlicher Gewinn pro Produkt: \antwortfeld[2.5cm]
\end{enumerate}

\section*{Etappe 7: Diagramm erstellen}
\begin{boxHinweis}
\textbf{Aufgabe.} Erstelle ein Säulendiagramm zum Gewinn je Artikel.
\end{boxHinweis}

\begin{boxK}
\textbf{Checkliste:}
\begin{itemize}
\item[$\square$] Ich habe die Artikelnamen \code{A3:A8} markiert.
\item[$\square$] Ich habe die Gewinne \code{G3:G8} markiert.
\item[$\square$] Ich habe ein Säulendiagramm eingefügt.
\item[$\square$] Mein Diagramm heißt \textbf{Gewinn je Artikel}.
\end{itemize}
\end{boxK}

\begin{boxTipp}
\textbf{Hilfe-Verweis:} Für das Diagramm brauchst du Bereiche. Wiederhole in der Excel-Hilfe, was \code{A3:A8} und \code{G3:G8} bedeuten.
\end{boxTipp}

\textbf{Kontrollfragen:}
\begin{enumerate}[label=\alph*)]
\item Welche Bereiche brauchst du? \linie[8cm]
\item Welcher Artikel hat die höchste Säule? \antwortfeld[3cm]
\end{enumerate}

\medskip
\section*{Etappe 8: Ergebnisse deuten}
\begin{boxHinweis}
\textbf{Aufgabe.} Schreibe eine kurze Auswertung. Nutze Zahlen aus deiner Tabelle.
\end{boxHinweis}

\textbf{Meine Auswertung:}\\[0.8cm]
\linie\\[0.8cm]
\linie\\[0.8cm]
\linie\\[0.8cm]
\linie

\begin{boxTipp}
\textbf{Hilfe-Verweis:} Prüfe vor deiner Begründung die Tabelle: Formelstart, Zellbezüge, SUMME und Gewinnspalte stehen in der Excel-Hilfe.
\end{boxTipp}

\textbf{Kontrollfragen:}
\begin{enumerate}[label=\alph*)]
\item Größter Gesamtgewinn bei: \antwortfeld[3cm]
\item Neben dem Verkaufspreis sind wichtig: \linie[7cm]
\end{enumerate}

\section*{Etappe 9: Sprinter -- Rabatt mit absolutem Zellbezug}
\begin{boxHinweis}
\textbf{Aufgabe.} Trage in B13 den Rabatt 10\% ein. Berechne den Gewinn nach Rabatt mit einer festen Rabattzelle.
\end{boxHinweis}

\begin{tabularx}{\textwidth}{|L{5cm}|Y|}
\hline
\rowcolor{gray!20}\textbf{Frage} & \textbf{Meine Antwort}\\\hline
Warum steht \texttt{\$B\$13} in der Formel? & \\\hline
Welche Formel steht in H3? & \\\hline
Gesamtgewinn nach Rabatt & \\\hline
Um wie viel Euro sinkt der Gewinn? & \\\hline
\end{tabularx}

\begin{boxTipp}
\textbf{Hilfe-Verweis:} Schaue in der Excel-Hilfe bei \textbf{Feste Zelle mit \$}. Dort steht, warum \texttt{\$B\$13} beim Herunterziehen gleich bleibt.
\end{boxTipp}

\section*{Abschluss-Check}
\begin{boxK}
\begin{itemize}
\item[$\square$] Ich kann eine Zelle über Spalte und Zeile finden.
\item[$\square$] Ich kann Messwerte korrekt eintragen.
\item[$\square$] Ich kann erklären, warum \code{=C3*D3} den Umsatz berechnet.
\item[$\square$] Ich kann \code{SUMME}, \code{MITTELWERT} und Bereiche wie \code{G3:G8} nutzen.
\item[$\square$] Ich nutze die Excel-Hilfe, wenn ich nicht weiterkomme.
\end{itemize}
\end{boxK}
